\section{引言}
全切片图像（WSIs）以高分辨率记录大范围组织形态，可用于肿瘤测量、组织组成分析与感兴趣区域选择。HistoSegNet和CAMEL等弱监督组织分割方法降低了标注需求，但预测的仍是预先定义的组织类别~\cite{chan2019histosegnet,xu2019camel}。复杂的组织边界与显著的尺度变化使密集像素标注成本高昂。此外，分割目标会随研究需求变化：同一张切片可能需要分别提取肿瘤上皮、间质或脂肪组织。固定的分割输出通道无法直接表达这种由用户定义的目标。

基于点击的交互方法与可提示基础模型提供了直接指定分割目标的方式~\cite{sofiiuk2022ritm,liu2023simpleclick,kirillov2023sam}。但将其应用于整张WSI仍面临三个挑战。首先，同一种组织可能分布于多个不连通区域，因此局部提示需要将信息传递至其邻域之外。其次，不同组织可能具有相似的局部外观，目标识别因而依赖周围的组织结构以及应予排除的示例。第三，组织识别需要较大范围的上下文，而小结构的轮廓描绘需要局部细节~\cite{tokunaga2019awmf,vanrijthoven2021hooknet,ho2021dmmn}。在整片尺度上，重复添加提示和重复进行图像计算进一步加重了这些问题。

因此，少量点击驱动的WSI分割需要在空间分离的区域中找回目标，同时限制目标缺失区域中的误检。当一个支持示例引导整张切片的预测时，目标覆盖与误检之间的权衡尤为关键。FewClick结合正负支持证据与密集概率迁移来处理这两方面需求。整块与分块特征保留组织结构及局部细节，多尺度区域上下文辅助目标识别，对应关系则将支持证据关联至远端查询位置。

图~\ref{fig:overview}概述了FewClick。共享骨干网络编码对齐的精细、中间和粗尺度视野，并融合整块与分块精细特征。区域池化与跨尺度消息组织上下文信息。正负支持点击引导初始预测，随后由粗到细的对应关系迁移预测的支持区域概率，细化查询logits。同一支持示例用于整片预测，无需追加人工提示。

本文的贡献包括三个方面。(1) 提出整片组织传播框架，仅需在一个支持图像块中提供少量正负点击，查询区域无需追加人工提示。(2) 结合多尺度整块／分块特征、正负支持证据与由粗到细的概率迁移，兼顾目标覆盖与非目标排除，并利用区域上下文组织空间信息。(3) 在共同的支持提示与新增适配预算下，基于15张WSI和11,000个远端开发任务比较六种方法，并通过匹配对照量化覆盖–误检权衡。FewClick在所比较的方法中取得最高的整片宏平均Dice（43.75\%），比SANSA高10.52个百分点。在远端开发队列上，其目标缺失FPA最低（7.75\%），同时保持最高的目标存在宏平均Dice（42.11\%）。
